\section*{Bab \ref{ch:g2:mixing-and-dissolving} --- Mencampur dan Melarutkan}

\begin{solution}{exo:g2:mixing-and-dissolving:1}
Gula dan garam \omterm{def:g2:mixing-and-dissolving:dissolve}{larut}. Pasir, kapur tulis, dan kerikil tidak \omterm{def:g2:mixing-and-dissolving:dissolve}{larut}.
\end{solution}

\begin{solution}{exo:g2:mixing-and-dissolving:2}
\omterm{def:g2:mixing-and-dissolving:dissolve}{Larutan} adalah cairan jernih yang diperoleh ketika suatu padatan \omterm{def:g2:mixing-and-dissolving:dissolve}{larut}
dalam air. Contoh: air gula, air asin (teh manis juga termasuk).
\end{solution}

\begin{solution}{exo:g2:mixing-and-dissolving:3}
Ya: gandum, kismis, kacang, dan biji-bijian disatukan. Setiap bagiannya
masih terlihat dan masih dapat dipungut.
\end{solution}

\begin{solution}{exo:g2:mixing-and-dissolving:4}
Tidak. Padatan yang \omterm{def:g2:mixing-and-dissolving:dissolve}{larut} membiarkan air tetap jernih tanpa apa-apa di
dasar; di sini airnya keruh dan ada lapisan yang mengendap.
\end{solution}

\begin{solution}{exo:g2:mixing-and-dissolving:5}
Air itu jernih (kita dapat melihat menembusnya), tetapi tidak tak
berwarna (warnanya merah muda). Air itu \omterm{def:g2:mixing-and-dissolving:dissolve}{larutan}: \omterm{def:g2:mixing-and-dissolving:dissolve}{larutan} boleh berwarna.
\end{solution}

\begin{solution}{exo:g2:mixing-and-dissolving:6}
Tehnya terasa manis: gula masih ada di dalamnya, tersebar di seluruh air
dalam bagian-bagian yang terlalu kecil untuk dilihat.
\end{solution}

\begin{solution}{exo:g2:mixing-and-dissolving:7}
Matahari mengeringkan airnya; garam yang tadinya \omterm{def:g2:mixing-and-dissolving:dissolve}{larut} dalam air laut
tertinggal.
\end{solution}

\begin{solution}{exo:g2:mixing-and-dissolving:8}
``\omterm{def:g2:mixing-and-dissolving:dissolve}{Larut}.'' Mencair adalah yang dilakukan es batu dengan sendirinya,
tanpa air; gula memerlukan air, dan gula itu masih ada di dalam teh.
\end{solution}

\begin{solution}{exo:g2:mixing-and-dissolving:9}
Gelas kimia berisi pasir keruh dan ada lapisan di dasarnya; air gula
jernih dan dasarnya kosong.
\end{solution}

\begin{solution}{exo:g2:mixing-and-dissolving:10}
Air keran bukan hanya air: ada sesuatu yang \omterm{def:g2:mixing-and-dissolving:dissolve}{larut} di dalamnya. Ketika
airnya mengering, padatan itu tertinggal sebagai cincin putih.
\end{solution}

\begin{solution}{exo:g2:mixing-and-dissolving:11}
Guru itu membiarkan beberapa tetes dari masing-masing gelas mengering
di atas cawan gelap yang bersih. Air murni tidak meninggalkan apa-apa;
air asin meninggalkan garam putih.
\end{solution}
